\documentclass[10pt,a4paper]{article}
\usepackage[utf8]{inputenc}
\usepackage[margin=0.85in]{geometry}
\usepackage{amsmath,amssymb}
\usepackage{booktabs}
\usepackage{tabularx}
\usepackage{array}
\usepackage{xcolor}
\usepackage{xurl}
\usepackage{hyperref}
\usepackage{enumitem}
\usepackage{titlesec}
\usepackage{tcolorbox}

\hypersetup{
    colorlinks=true,
    linkcolor=blue!70!black,
    citecolor=blue!70!black,
    urlcolor=blue!70!black
}

\titlespacing{\section}{0pt}{10pt}{4pt}
\titlespacing{\subsection}{0pt}{8pt}{3pt}
\titlespacing{\subsubsection}{0pt}{6pt}{2pt}

\tcbset{
    boxrule=0.5pt,
    arc=2pt,
    left=6pt,
    right=6pt,
    top=4pt,
    bottom=4pt
}

\newtcolorbox{reviewercomment}{
    colback=gray!8,
    colframe=gray!40,
    fonttitle=\bfseries\small,
    title=Reviewer Comment
}

\begin{document}

\begin{center}
    {\LARGE \textbf{Author Response to Reviewers}}\\[1ex]
    {\large \textbf{Paper ID: 1021}}\\[0.5ex]
    \textbf{Title:} Empirical Analysis of Optimization Phase Ordering Trade-offs in a Custom Compiler Using Micro-Benchmarks\\[0.5ex]
    \textbf{Track:} Computing --- COSMIC 2026\\[0.5ex]
    \textit{Authors:} Aniket De, Kushal S M, Manas S, Vivek Neeralagi, Suhas A Bhyratae, Harisha
\end{center}

\vspace{0.5em}
\hrule
\vspace{1em}

\section*{Executive Summary \& Overview of Revisions}

We sincerely thank the Track Chair, Associate Editors, and Reviewers for their thoughtful, rigorous, and constructive feedback on our submission. We are greatly encouraged by the positive reception---including recommendations of \textbf{Accept}, \textbf{Accept}, and \textbf{Weak Accept}---and by the reviewers' recognition of our modular compiler architecture, the utility of our independently reorderable pass framework, and our clear visualization of pass-ordering dynamics.

The reviewers highlighted crucial areas where the initial draft needed expansion, specifically:
\begin{enumerate}[itemsep=1pt,parsep=0pt,topsep=2pt]
    \item \textbf{Generalizability beyond a single benchmark:} Evaluating across the full promised suite of MiBench and PolyBench kernels rather than focusing primarily on \texttt{factorial}.
    \item \textbf{Explicit specification of the execution cost model:} Formally defining the per-instruction cycle weights and introducing model-free validation.
    \item \textbf{Rigorous statistical substantiation:} Formally reporting statistical tests, test statistics, $p$-values, and effect sizes.
    \item \textbf{Clarification of pairwise interaction matrices and units:} Explaining the derivation, sample space, and interpretation of Table~IV.
    \item \textbf{Baseline transparency:} Providing complete source code and intermediate representation (IR) line-by-line accounting for the baseline measurements.
    \item \textbf{Clarification of the speed--size trade-off and novelty scope:} Explaining the Pareto characteristics of code-shrinking passes and explicitly delimiting our contribution.
\end{enumerate}

In response, we undertook an exhaustive overhaul of the experimental evaluation and manuscript:
\begin{itemize}[itemsep=2pt,parsep=0pt,topsep=2pt]
    \item \textbf{From 1 Benchmark to 15 Benchmarks (10,815 Executions):} We expanded our evaluation to \textbf{fifteen benchmark programs} (8 scalar kernels adapted from MiBench and 7 array/matrix kernels adapted from PolyBench/C). We evaluated \textbf{all $6! = 720$ phase orderings} across all 15 programs, generating \textbf{10,800 optimized configurations} (plus 15 baselines). Every single optimized binary was dynamically executed and verified against baseline output by an automated test harness to guarantee semantic equivalence.
    \item \textbf{Formal Instruction Cost Model \& Model-Free Dynamic Validation:} We added Section~III-C and Table~II, explicitly detailing the per-opcode cycle latency for all 21 3AC instructions based on an in-order ARM Cortex-M architecture. Furthermore, we integrated an IR execution interpreter to measure \textbf{Dynamic Instruction Count} on fixed inputs as an independent, model-free runtime metric ($\rho = 0.88$ correlation with cycle estimates).
    \item \textbf{Complete Statistical Validation:} We added Section~IV-D reporting non-parametric hypothesis tests across the 15 program blocks: Friedman tests for code size ($\chi^2(23) = 304.6, p = 5.3 \times 10^{-51}, W = 0.88$), estimated cycles ($p = 3.0 \times 10^{-49}, W = 0.86$), and dynamic instructions ($p = 8.5 \times 10^{-51}, W = 0.88$), followed by an exact Wilcoxon signed-rank post-hoc test ($p = 3.1 \times 10^{-5}, r = 1.00$).
    \item \textbf{Formalization of Pairwise Interaction Matrix:} Section~IV-B and Table~IV define the exact derivation across all 720 orderings, specifying units (mean IR instructions saved/lost) and explaining marginal vs.\ unconfounded effects.
    \item \textbf{Baseline Line-by-Line Breakdown:} Section~IV-A and Listing~1 provide the complete source code, unoptimized IR breakdown (40 instructions, 81 cycles), and optimized IR (16 instructions, 55 cycles) for \texttt{factorial}.
    \item \textbf{Single-Command Full Reproducibility:} The entire compiler, benchmark suite, and figure-generation scripts are open-sourced at \url{https://github.com/VivekNeer/major-project-compiler-design}. All 10,815 compilations and tests can be reproduced with a single command (\texttt{python -m compiler.main --benchmark-all}) in $\sim$23 seconds.
\end{itemize}

\vspace{1em}
\section*{Detailed Response to Reviewer \#1}

\begin{reviewercomment}
\textbf{Comment 1 (Single benchmark):} Single benchmark. All reported results (Tables~II--V, Figs.~2--5) come from one program (factorial), despite the stated objective of validating against MiBench and PolyBench benchmarks. This significantly limits the generalizability of the conclusions.
\end{reviewercomment}

\noindent \textbf{Response:} We completely agree with the reviewer. In the revised manuscript, we have addressed this comprehensively by expanding our empirical evaluation to \textbf{fifteen benchmark programs} (Table~V and Section~IV-C):
\begin{itemize}[itemsep=1pt,parsep=0pt,topsep=2pt]
    \item \textbf{8 MiBench-derived scalar kernels:} \texttt{factorial}, \texttt{isqrt}, \texttt{bitcount}, \texttt{fibonacci}, \texttt{collatz}, \texttt{sha\_mix}, \texttt{power}, and \texttt{gcd}.
    \item \textbf{7 PolyBench/C-derived array and stencil kernels:} \texttt{jacobi1d}, \texttt{jacobi2d}, \texttt{gemm}, \texttt{2mm}, \texttt{gesummv}, \texttt{atax}, and \texttt{bicg}.
\end{itemize}
Rather than testing a subset of sequences, we executed \textbf{all $6! = 720$ permutations across all 15 benchmarks}, totaling \textbf{10,800 optimized configurations} (10,815 total runs with baselines).
\begin{enumerate}[itemsep=1pt,parsep=0pt,topsep=2pt]
    \item \textbf{Per-Program Reductions (Table~V):} The best ordering achieves a geometric-mean reduction of \textbf{35.7\% in static code size} (ranging from 22.3\% on \texttt{bicg} to 60.0\% on \texttt{factorial}), \textbf{20.5\% in estimated cycles} (11.9\% to 33.9\%), and \textbf{23.2\% in dynamic instruction count} (12.9\% to 34.6\%).
    \item \textbf{Distribution across Orderings (Figure~3):} Box plots illustrate the spread across all 720 orderings for each of the 15 kernels. On 13 of the 15 programs, the worst orderings produce 0\% reduction, confirming that poor phase ordering completely nullifies optimization opportunities.
    \item \textbf{Universality of Equivalence Classes (Table~VI):} Across all 15 programs, the 720 orderings collapse into exactly \textbf{nine equivalence classes} determined by the relative order of the four size-relevant passes ($\text{CF}, \text{CSE}, \text{CP}, \text{DCE}$). The optimal relative order ($\text{CF} \to \text{CSE} \to \text{CP} \to \text{DCE}$) is identical across all 15 benchmarks.
\end{enumerate}
\texttt{factorial} is now retained strictly as an introductory worked example in Section~IV-A to allow readers to verify the line-by-line mechanics before presenting the cross-benchmark results in Sections~IV-C, IV-D, and IV-E.

\vspace{0.8em}
\begin{reviewercomment}
\textbf{Comment 2 (Undefined cost model):} ``Estimated execution cycles'' is a primary metric but the weighted instruction-cost model behind it is never specified (what does each instruction type cost?), making the cycle results unverifiable.
\end{reviewercomment}

\noindent \textbf{Response:} We thank the reviewer for highlighting this omission. In the revised paper, we added \textbf{Section~III-C (``Metrics and Cost Model'')} and \textbf{Table~II (``Per-instruction cost model for the estimated-cycle metric'')}, which fully specifies the per-opcode cycle latency model:
\begin{itemize}[itemsep=1pt,parsep=0pt,topsep=2pt]
    \item Single-cycle ALU operations (\texttt{ADD}, \texttt{SUB}, \texttt{NEG}, comparisons, \texttt{AND}, \texttt{OR}, \texttt{NOT}, \texttt{COPY}, \texttt{LOAD\_CONST}, \texttt{PARAM}): \textbf{1 cycle}.
    \item Multi-cycle multiplier (\texttt{MUL}): \textbf{3 cycles}.
    \item Iterative divider/modulo (\texttt{DIV}, \texttt{MOD}): \textbf{12 cycles}.
    \item Control flow branches \& jumps (\texttt{JUMP}, \texttt{JUMP\_IF\_TRUE}, \texttt{JUMP\_IF\_FALSE}): \textbf{2 cycles}.
    \item Memory loads \& stores (\texttt{ARR\_LOAD}, \texttt{ARR\_STORE}, \texttt{GLOBAL\_LOAD}, \texttt{GLOBAL\_STORE}): \textbf{2 cycles}.
    \item Function calls, returns, prints: \texttt{CALL} = \textbf{5}, \texttt{RETURN} = \textbf{3}, \texttt{PRINT} = \textbf{10 cycles}.
    \item Non-executable structural markers (\texttt{LABEL}, \texttt{NOP}, \texttt{FUNC\_BEGIN}, etc.): \textbf{0 cycles}.
\end{itemize}
\textbf{Model Rationale \& Model-Free Validation:}
\begin{enumerate}[itemsep=1pt,parsep=0pt,topsep=2pt]
    \item \textit{Target Architecture:} The weights reflect the typical in-order execution latencies of an ARM Cortex-M embedded core, directly aligned with embedded workloads represented by MiBench.
    \item \textit{Model-Free Validation:} To ensure our conclusions do not depend on synthetic latency weights, we introduced \textbf{Dynamic Instruction Count}---the exact count of 3AC instructions executed at runtime by our interpreter on fixed inputs. As shown in Section~IV-E and Table~VI, the static cycle metric and the dynamic instruction count are strongly correlated ($\rho = 0.88$), confirming that the cycle model reliably tracks execution cost without distorting ordering ranks.
\end{enumerate}

\vspace{0.8em}
\begin{reviewercomment}
\textbf{Comment 3 (Unsubstantiated statistical claims):} Section~IV.D states that statistical tests assessed significance, but no test, p-value, or effect size is reported anywhere in the results section.
\end{reviewercomment}

\noindent \textbf{Response:} We apologize for the omission of statistical reporting in the initial submission. In the revised paper, we added \textbf{Section~IV-D (``Statistical Significance'')}, presenting a complete non-parametric statistical evaluation treating the fifteen benchmark programs as independent blocks/replications:
\begin{enumerate}[itemsep=1pt,parsep=0pt,topsep=2pt]
    \item \textbf{Friedman Test (Repeated-Measures Non-Parametric ANOVA):} We tested whether phase ordering significantly influences compiler performance across the 15 benchmark programs, evaluating the 24 relative orderings of the four size-relevant passes ($\text{CF}, \text{CSE}, \text{CP}, \text{DCE}$):
    \begin{itemize}[itemsep=1pt,parsep=0pt]
        \item \textbf{Static Code Size:} $\chi^2(23) = 304.6,\; p = 5.3 \times 10^{-51},\; \text{Kendall's } W = 0.88$
        \item \textbf{Estimated Cycles:} $\chi^2(23) = 296.0,\; p = 3.0 \times 10^{-49},\; \text{Kendall's } W = 0.86$
        \item \textbf{Dynamic Instruction Count:} $\chi^2(23) = 303.6,\; p = 8.5 \times 10^{-51},\; \text{Kendall's } W = 0.88$
    \end{itemize}
    \textit{Interpretation:} The null hypothesis of equal performance across orderings is rejected at $p < 10^{-48}$. Kendall's $W \in [0.86, 0.88]$ indicates exceptionally strong concordance across all fifteen programs, formally proving that the ordering hierarchy is consistent across diverse programs.
    \item \textbf{Post-Hoc Pairwise Comparison (Wilcoxon Signed-Rank Test):} An exact Wilcoxon signed-rank test comparing Class 1 (optimal) against Class 9 (worst) across the 15 paired programs yielded:
    \[ W = 0,\quad p = 3.1 \times 10^{-5} \text{ (one-sided exact)},\quad \text{rank-biserial } r = 1.00 \]
    Class 1 achieved superior reduction over Class 9 on every single program ($15/15$), yielding the minimum possible $p$-value for $n = 15$ with maximal effect size ($r = 1.00$) and a median paired improvement of 34\% of baseline code size.
    \item \textbf{Spearman Rank Correlation ($\rho$):} Reported across all 10,800 optimized configurations: $\rho = 0.79$ between code size and estimated cycles; $\rho = 0.88$ between code size and dynamic instruction count (median within-program $\rho = 1.00$).
\end{enumerate}

\vspace{0.8em}
\begin{reviewercomment}
\textbf{Comment 4 (Table IV underspecified):} Table~IV is underspecified. The pairwise pass-interaction values are not explained in terms of derivation or units; with only 6 evaluated sequences it is unclear how a full 6$\times$6 interaction matrix was populated.
\end{reviewercomment}

\noindent \textbf{Response:} We thank the reviewer for identifying this ambiguity. In the revised paper, we clarified this in \textbf{Section~IV-B (``Pairwise Pass Interactions'')} and the caption of \textbf{Table~IV}:
\begin{enumerate}[itemsep=1pt,parsep=0pt,topsep=2pt]
    \item \textbf{Derivation \& Sample Size:} The pairwise matrix is \textbf{not} computed from the 6 illustrative sequences; it is computed across \textbf{all 720 full permutations} on \texttt{factorial}.
    \item \textbf{Mathematical Definition \& Units:} For each ordered pair of passes $(A, B)$, the entry $M_{AB}$ is defined as:
    \[ M_{AB} = \overline{\text{Size}}_{A \prec B} - \overline{\text{Size}}_{B \prec A} \]
    where $\overline{\text{Size}}_{A \prec B}$ is the mean static code size (in \textbf{IR instructions}) across the 360 orderings in which pass $A$ precedes pass $B$, and $\overline{\text{Size}}_{B \prec A}$ is the mean over the 360 orderings in which $B$ precedes $A$.
    \item \textbf{Interpretation:} Negative values mean placing pass $A$ before pass $B$ produces fewer instructions (e.g., $M_{\text{CF}, \text{DCE}} = -10.7$, indicating that running Constant Folding before Dead Code Elimination saves $10.7$ instructions on average). The matrix is anti-symmetric ($M_{AB} = -M_{BA}$) with diagonal elements equal to 0 by construction.
    \item \textbf{Marginality Caveat:} We explicitly discuss in Section~IV-B that this pairwise matrix is a marginal statistic averaging over the positions of the other four passes. We explain why passes like Algebraic Simplification (AS) and Strength Reduction (SR) show small non-zero values ($+2.9$): in the permutation space, having AS before CF correlates with DCE appearing later in the pipeline. We guide the reader to the unconfounded Equivalence Class analysis (Section~IV-C / Table~VI) for exact causal attribution.
\end{enumerate}

\vspace{0.8em}
\begin{reviewercomment}
\textbf{Comment 5 (Baseline plausibility):} Baseline plausibility. A baseline instruction count of 40 for an unoptimized factorial routine is not shown in source or IR form, making it hard for readers to verify the reported 60\%/32\% reduction figures.
\end{reviewercomment}

\noindent \textbf{Response:} We have made the entire derivation transparent in \textbf{Section~IV-A} and \textbf{Listing~1}:
\begin{enumerate}[itemsep=1pt,parsep=0pt,topsep=2pt]
    \item \textbf{Source Code Included (Listing~1):} The full source code of the \texttt{factorial} program (containing both \texttt{factorial(n)} and \texttt{main()}) is shown in Listing~1. It explicitly contains constant folding opportunities (\texttt{base = 2 + 3}, \texttt{extra = 10 - 5}), dead code assignments (\texttt{unused = 42 * 2}, \texttt{also\_unused = unused + 100}), and repeated subexpressions (\texttt{sum1 = base + extra}, \texttt{sum2 = base + extra}).
    \item \textbf{Instruction Breakdown of the 40 Baseline Instructions:} Because the unoptimized 3AC lowering places every constant and intermediate into a dedicated temporary (e.g., \texttt{t6 = 2; t7 = 3; t8 = t6 + t7; base = t8}), the 40 countable instructions consist of: 11 constant loads (\texttt{LOAD\_CONST}), 11 variable copies (\texttt{COPY}), 5 additions (\texttt{ADD}), 2 multiplications (\texttt{MUL}), 1 subtraction (\texttt{SUB}), 1 comparison (\texttt{LTE}), 2 jumps (\texttt{JUMP}, \texttt{JUMP\_IF\_FALSE}), 1 parameter setup (\texttt{PARAM}), 1 function call (\texttt{CALL}), 2 returns (\texttt{RETURN}), 3 prints (\texttt{PRINT}) (13 instructions in \texttt{factorial()}, 27 instructions in \texttt{main()}).
    \item \textbf{Cycle Accounting:} Applying Table~II:
    \[ \text{Cycles} = (30 \times 1) + (2 \times 3) + (2 \times 2) + 5 + (2 \times 3) + (3 \times 10) = 30 + 6 + 4 + 5 + 6 + 30 = 81 \text{ cycles.} \]
    \item \textbf{Optimized Code:} The optimal sequence reduces the loop in \texttt{factorial} to 10 instructions and \texttt{main()} to 6 instructions (\texttt{param 5; t12 = call factorial; print t12; print 10; print 10; return 0}), yielding \textbf{16 instructions} and \textbf{55 cycles}.
    \begin{itemize}[itemsep=1pt,parsep=0pt]
        \item Size reduction: $(40 - 16) / 40 = 60.0\%$.
        \item Cycle reduction: $(81 - 55) / 81 = 32.1\%$.
    \end{itemize}
\end{enumerate}

\vspace{0.8em}
\begin{reviewercomment}
\textbf{Comment 6 (Minor labeling and normalization):} Minor: inconsistent table/sequence labeling (Table~II sequences aren't explicitly numbered ``Seq~1--6'' as referenced later in Tables~III/V); normalization method for Table~V is not stated explicitly.
\end{reviewercomment}

\noindent \textbf{Response:} Both minor issues have been resolved:
\begin{itemize}[itemsep=1pt,parsep=0pt,topsep=2pt]
    \item In \textbf{Table~III}, all evaluated sequence rows are explicitly and consistently labeled \texttt{Base}, \texttt{Seq 1}, \texttt{Seq 2}, \texttt{Seq 3}, \texttt{Seq 4}, \texttt{Seq 5}, and \texttt{Seq 6}.
    \item In \textbf{Section~III-C}, the normalization formula is explicitly stated:
    Percentage reduction: $100 \times (1 - m_{\mathrm{opt}} / m_{\mathrm{base}})$;
    Normalized ratio: $m_{\mathrm{opt}} / m_{\mathrm{base}}$ (where 1.0 represents baseline; smaller is better).
    Cross-program averages (Table~VI) explicitly use the \textbf{geometric mean} of normalized ratios, following Fleming \& Wallace (1986).
\end{itemize}

\vspace{1em}
\section*{Detailed Response to Reviewer \#2}

\begin{reviewercomment}
\textbf{Comment 1 (Evaluation Scope, Cost Model, and Statistical Tests):} The evaluation relies mainly on one factorial benchmark, unvalidated cycle estimates, and statistical tests whose design and results are not reported, limiting support for general claims about phase ordering.
\end{reviewercomment}

\noindent \textbf{Response:} We thank the reviewer for these critical observations. We have resolved all three concerns in depth:
\begin{enumerate}[itemsep=1pt,parsep=0pt,topsep=2pt]
    \item \textbf{15 Diverse Benchmarks \& 10,815 Runs:} Evaluated 8 scalar kernels from MiBench and 7 array/stencil kernels from PolyBench/C across all 720 orderings (Table~V, Fig.~3, Table~VI).
    \item \textbf{Explicit Cost Model \& Model-Free Dynamic Validation:} Defined the full ARM Cortex-M latency model (Table~II) and validated it against dynamic instruction counts executed by an IR interpreter ($\rho = 0.88$). In addition, in Section~IV-F, we lowered baseline and optimized IR to physical RV32IM RISC-V assembly, confirming that the optimizations translate to hardware code generation with a 26.0\% geometric-mean reduction in assembly instructions.
    \item \textbf{Rigorous Non-Parametric Statistics:} Formally reported the Friedman test ($\chi^2 = 304.6, p < 10^{-48}, W = 0.88$) and Wilcoxon signed-rank test ($p = 3.1 \times 10^{-5}, r = 1.00$) in Section~IV-D.
\end{enumerate}

\vspace{0.8em}
\begin{reviewercomment}
\textbf{Comment 2 (Speed--Size Trade-Off, Pass Descriptions, and Contribution):} The paper should clarify the claimed speed--size trade-off, describe the added passes and evaluation in reproducible detail, and distinguish its contribution from established compiler techniques.
\end{reviewercomment}

\noindent \textbf{Response:} We have addressed each of these three points:
\begin{enumerate}[itemsep=2pt,parsep=0pt,topsep=2pt]
    \item \textbf{Clarification of Speed--Size Trade-Off (Section~IV-E):} We clarified why the Pareto front between code size and execution performance is narrow in this study. The six passes evaluated here (CF, DCE, CSE, CP, SR, AS) are \textbf{strictly code-shrinking transformations}. Unlike code-expanding optimizations (such as loop unrolling, function inlining, or loop tiling) that trade larger binary footprint for fewer runtime branches or higher ILP, code-shrinking passes reduce execution cycles proportionally by eliminating redundant operations. Consequently, Class 1 simultaneously minimizes code size (0.643), estimated cycles (0.795), and dynamic instructions (0.768) on all 15 programs ($\rho = 0.79 \text{ to } 0.88$). We explicitly contrast this with multi-objective frameworks like POSET-RL on LLVM, noting that the absence of a size-vs-speed conflict is an intrinsic property of this pass set.
    \item \textbf{Detailed Pass Descriptions \& Reproducibility (Section~III-B \& Section~IV):}
    Section~III-B details the preconditions, liveness analysis, available-expression computation, and rewriting rules for each pass. All passes are implemented as block-local, side-effect-free transformations that preserve semantic correctness under any ordering. The complete pure-Python compiler, test suite (161 unit tests), and benchmarking pipeline are open-sourced at \url{https://github.com/VivekNeer/major-project-compiler-design}. The entire experiment (10,815 compile-and-execute cycles) can be regenerated with a single command (\texttt{python -m compiler.main --benchmark-all}) in $\sim$23 seconds.
    \item \textbf{Distinguishing Our Contribution:} We explicitly state in Section~I and Section~VI that our contribution is \textbf{not} the invention of new transformation passes (which are textbook techniques), but rather:
    \begin{itemize}[itemsep=1pt,parsep=0pt]
        \item \textit{Orthogonal Compiler Architecture:} Designing an infrastructure where passes are strictly isolated and side-effect free, allowing all $6! = 720$ orderings to run without compiler crashes or semantic divergence.
        \item \textit{Exhaustive Empirical Structure:} Discovering that the 720 orderings collapse into exactly 9 equivalence classes governed by 4 fundamental ordering rules.
        \item \textit{Universality Proof:} Proving empirically and statistically that the relative ordering $\text{CF} \to \text{CSE} \to \text{CP} \to \text{DCE}$ is universally optimal across both scalar and polyhedral benchmark suites.
    \end{itemize}
\end{enumerate}

\vspace{1em}
\section*{Detailed Response to Reviewer \#4}

\begin{reviewercomment}
\textbf{Assessment:} The paper demonstrates a clean and modular compiler architecture, as illustrated in Fig.~1, where each optimization pass is implemented as an independently reorderable component, providing a useful framework for controlled experimentation and analysis of pass interactions. The use of Three-Address Code (3AC) as the intermediate representation is appropriate for supporting classical data-flow optimizations such as liveness analysis for Dead Code Elimination and available-expression analysis for Common Subexpression Elimination. The study also makes effective use of visualizations, including scatter plots, heatmaps, and normalized bar charts, to clearly communicate the trade-offs between code size and estimated execution cycles across different optimization sequences. \textit{(Overall Decision: Accept)}
\end{reviewercomment}

\noindent \textbf{Response:} We warmly thank Reviewer \#4 for the positive evaluation and for recommending \textbf{Accept}. We appreciate the reviewer's recognition of our modular architecture, the suitability of 3AC IR for classical data-flow analyses, and our visual presentation.

\textit{(Note: In the review submission portal, the text under ``Weaknesses'' was an exact duplicate of ``Strengths''.)} We have preserved all the architectural clarity and visual representations praised by the reviewer, while expanding the visualizations to include:
\begin{itemize}[itemsep=1pt,parsep=0pt,topsep=2pt]
    \item Figure~2: The $6 \times 6$ pairwise pass-interaction heatmap with explicit units.
    \item Figure~3: Full-width box plot distributions across all 720 orderings for all 15 benchmark programs.
\end{itemize}

\vspace{1em}
\section*{Summary Mapping of Review Comments to Final Manuscript}

\begin{table}[htbp]
\centering
\small
\renewcommand{\arraystretch}{1.1}
\begin{tabularx}{\textwidth}{@{}llXl@{}}
\toprule
\textbf{Reviewer \& Item} & \textbf{Topic} & \textbf{Summary of Resolution} & \textbf{Location in Paper} \\
\midrule
R1.1, R2.1 & Benchmark scope & Expanded from 1 to 15 benchmarks (MiBench + PolyBench/C), 720 orderings, 10,800 runs & Section~IV-C, Table~V, Fig.~3 \\
R1.2, R2.1 & Cost model & Complete opcode cost table (Table~II) + model-free dynamic instruction counts & Section~III-C, Table~II, Sec.~IV-E \\
R1.3, R2.1 & Statistical tests & Friedman test ($p < 10^{-48}$), Wilcoxon test ($p = 3.1 \times 10^{-5}$), Kendall's $W$ & Section~IV-D \\
R1.4 & Table~IV specification & Formal mathematical definition, units (IR instructions), and marginality caveats & Section~IV-B, Table~IV, Fig.~2 \\
R1.5 & Baseline plausibility & Full source code listing (Listing~1), line-by-line IR opcode and cycle breakdown & Section~IV-A, Listing~1 \\
R1.6 & Labels \& normalization & Standardized labels (\texttt{Base}, \texttt{Seq 1--6}), geometric-mean normalization stated & Section~III-C, Table~III, Table~VI \\
R2.2 & Speed--size trade-off & Explained absence of Pareto conflict for code-shrinking passes & Section~IV-E \\
R2.2 & Reproducibility & Formal pass specifications and 1-command reproducible open-source release & Section~III-B, Sec.~IV, Sec.~VI \\
R2.2 & Novelty scope & Clarified focus on exhaustive measurement, equivalence classes, and architecture & Section~I, Section~VI \\
R4 & Overall architecture & Preserved strengths; expanded cross-program plots and equivalence class tables & Section~III-A, Sec.~IV-C, Figs.~1--3 \\
\bottomrule
\end{tabularx}
\end{table}

\end{document}
